\chapter{Coste de Oportunidad}
\section{Estructura de capital}
La corriente de flujos de efectivo que producen los activos de una empresa son su recurso básico.

Cuando la empresa se financia totalmente con acciones comunes, todos esos flujos de efectivo pertenecen a los accionistas.

Pero cuando la empresa emite títulos de deuda y de capital, los flujos de efectivo se dividen en dos corriente, una corriente relativamente segura que va a los tenedores de deuda y una corriente más riesgosa que va a los accionistas.

La combinación del financiamiento a través de deuda y de capital se conoce como estructura de capital de una empresa.

La estructura de capital no solo corresponde a ``deuda o capital''. Existen muchos tipos de deuda, y por lo menos dos tipos de capital (ordinario y preferente), además de los híbridos como, por ejemplo, los bonos convertibles, leasing.

La empresa puede emitir docenas de títulos en innumerables combinaciones. Así, se trata de maximizar el valor de mercado de la empresa.

\section{Deuda / Capital}

\section{Coste de oportunidad}

¿Cuál es la rentabilidad mínima que debo obtener por las acciones? Es una pregunta muy fre cuente que se hacen los inversionistas.

El dinero que financia una empresa proviene de diversas fuentes de financiación; pero ¿cuánto cuesta conseguir esos fondos?

En finanzas existe un principio básico que es que el dinero tiene valor en el tiempo.

El coste de oportunidad es la tasa mínima exigida por los accionistas para mantener su inversión en la empresa. Es aquella que el inversionista deja de ganar por tener el dinero invertido en la empresa, se puede expresar en términos de rentabilidad porque representa la rentabilidad que ofrece el mercado por otras inversiones de igual riesgo.

Es el beneficio que podría haber obtenido si hubiera elegido la alternativa más rentable --> comparar lo elegido frente a la mejor opción

\section{Coste del capital}
Es un error muy común creer que la tasa de descuento que se utiliza para analizar inversiones (VAN, TIR) es la tasa de interés.

El coste de oportunidad depende del riesgo del negocio, entonces: ¿Cuál es la tasa que refleja el riesgo de los flujos de caja?

\subsection{Modelo CAPM}
El CAPM analiza la relación entre el riesgo y el rendimiento esperando de un activo. 

Deriva del modelo de varianza de Markowitz, el cual 